\section{Dynamic Memory Allocation}  

\begin{definition} 
\textbf{Static vs Automatic Memory Allocation:} Global variables are statically allocated when the program is loaded and deallocated when it exits, while function variables are automatically allocated on function call and deallocated on return.
\end{definition}  

\begin{remark} 
Note that because automatically allocated memory is created on function calls, it is stored on the stack.
\end{remark}  

\begin{remark} 
Often we need a larger amount of memory space that persists across multiple function calls without lasting the entire program lifetime.
\end{remark}  

\begin{definition} 
\textbf{Dynamic Memory Allocation:} We can use allocation functions (such as \texttt{malloc}) to explicitly request dynamic memory blocks on the heap.
\end{definition}  

\begin{codeexample}
char *a = (char *) malloc(10 * sizeof(char)); // Allocates uninitialized memory for 10 characters
char *zeros = (char *) calloc(10, sizeof(char)); // Allocates zero-initialized memory for 10 characters
free(a); // Frees allocated memory block

zeros = (char *) realloc(zeros, 20 * sizeof(char)); // Resizes allocated memory block to 20 characters
\end{codeexample}  

\subsection{Data Structures}

\begin{definition} 
\textbf{Structured Data} is a custom type containing multiple data fields of different types. It is defined using the \texttt{struct} keyword.
\end{definition}  

\begin{codeexample}
struct myData { 
    int dataField; 
    char dataName;
};
\end{codeexample}  

\begin{definition} 
\textbf{Unions} overlay multiple member fields in a single shared memory location. Only one member value should be stored/interpreted at a time, and C performs no runtime tracking of the active member type.
\end{definition}  

\begin{codeexample}
union currentTime { 
    char stringName[10]; 
    int intTime; 
    long longTime;
};
\end{codeexample}  

\begin{remark} 
We can use \texttt{typedef} to define type aliases. Mixing struct tags and typedef aliases with the same name can cause confusion.
\end{remark}  

\begin{remark} 
Tags (e.g., \texttt{struct tag \{...\}}) belong to a separate namespace from typenames; consequently, a typename and a tag can share the same identifier.
\end{remark}  

\begin{remark} 
Structs and pointers allow building \textbf{dynamic data structures}. Using \texttt{void} pointers allows elements to be generic.
\end{remark}  

\begin{codeexample}
struct listNode { 
    void *value; 
    struct listNode *next;
};  

struct listNode *myList = NULL;
\end{codeexample}